\section{Ferrari's quartic formula}
\label{sec:quartic}

\subsection{Depressing the quartic}

Let
\[
 f(X)=aX^4+bX^3+cX^2+dX+e,
 \qquad a\ne0.
\]
Put
\[
 A=\frac ba,\quad B=\frac ca,\quad C=\frac da,\quad D=\frac ea
\]
and translate
\[
 X=Y-\frac A4.
\]
A direct expansion gives
\[
 \frac1a f\left(Y-\frac A4\right)=Y^4+pY^2+qY+r,
 \tag{F1}\label{eq:quartic-depress}
\]
where
\begin{align}
 p&=B-\frac{3A^2}{8},\label{eq:quartic-p}\\
 q&=C-\frac{AB}{2}+\frac{A^3}{8},\label{eq:quartic-q}\\
 r&=D-\frac{AC}{4}+\frac{A^2B}{16}-\frac{3A^4}{256}.
 \label{eq:quartic-r}
\end{align}
We therefore solve
\[
 h(Y)=Y^4+pY^2+qY+r.
\]

\subsection{The biquadratic branch}
\label{sec:quartic-biquadratic}

When $q=0$, no cubic resolvent is needed.  Apply the quadratic formula to
\[
 Z^2+pZ+r=0.
\]
Choose $\delta,z_+,z_-$ with
\[
 \delta^2=p^2-4r,
 \qquad
 z_+=\frac{-p+\delta}{2},
 \qquad
 z_-=\frac{-p-\delta}{2}.
\]
Then
\[
 Z^2+pZ+r=(Z-z_+)(Z-z_-).
\]
Choose square-root witnesses $\eta_+^2=z_+$ and $\eta_-^2=z_-$.  The exact
factorization is
\[
 Y^4+pY^2+r
 =(Y-\eta_+)(Y+\eta_+)(Y-\eta_-)(Y+\eta_-).
 \tag{F2}\label{eq:biquadratic-factor}
\]
This branch is entirely constructive once the equation $q=0$ is supplied; it
does not test equality of a newly introduced radical.

\subsection{Ferrari's factorization equations when \texorpdfstring{$q\ne0$}{q nonzero}}

Assume now that $q\ne0$.  Seek a factorization into two monic quadratics of
opposite linear terms:
\[
 h(Y)=\bigl(Y^2-sY+(m-t)\bigr)
       \bigl(Y^2+sY+(m+t)\bigr).
 \tag{F3}\label{eq:ferrari-target}
\]
Multiplication shows that~\eqref{eq:ferrari-target} follows from
\[
 s^2=2m-p,
 \qquad
 2st=-q,
 \qquad
 t^2=m^2-r.
 \tag{F4}\label{eq:ferrari-parameters}
\]
Indeed, the $Y^2$ coefficient is $2m-s^2=p$, the $Y$ coefficient is
$-2st=q$, and the constant term is $m^2-t^2=r$.

Eliminating $s$ and $t$ gives the cubic resolvent
\[
 q^2-4(2m-p)(m^2-r)=0.
 \tag{F5}\label{eq:ferrari-resolvent}
\]
Equivalently,
\[
 m^3-\frac p2m^2-rm+\frac{pr}{2}-\frac{q^2}{8}=0.
 \tag{F6}\label{eq:ferrari-resolvent-monic}
\]
Thus $m$ is obtained by the cubic construction of
Section~\ref{sec:cubic}.

Choose a resolvent root $m$ and a square-root witness
\[
 s^2=2m-p.
 \tag{F7}\label{eq:s-square}
\]
The hypothesis $q\ne0$ and~\eqref{eq:ferrari-resolvent} imply
$2m-p\ne0$: otherwise the right-hand product in
\eqref{eq:ferrari-resolvent} would vanish and hence $q^2=0$, forcing $q=0$.
Consequently $s\ne0$.  We may therefore define
\[
 t=-\frac{q}{2s}.
 \tag{F8}\label{eq:t-quotient}
\]
Then $2st=-q$.  The resolvent equation and~\eqref{eq:s-square} give
\[
 q^2=4(2m-p)(m^2-r)=4s^2(m^2-r),
\]
so
\[
 t^2=\frac{q^2}{4s^2}=m^2-r.
\]
All three equations in~\eqref{eq:ferrari-parameters} hold, and every division
has an explicit nonzero proof.

\subsection{Solving the two quadratics}

Choose square-root witnesses
\[
 \rho^2=s^2-4(m-t),
 \qquad
 \sigma^2=s^2-4(m+t).
 \tag{F9}\label{eq:rho-sigma}
\]
The first factor in~\eqref{eq:ferrari-target} has roots
\[
 y_0=\frac{s+\rho}{2},
 \qquad
 y_1=\frac{s-\rho}{2},
 \tag{F10}\label{eq:quartic-first-pair}
\]
and the second has roots
\[
 y_2=\frac{-s+\sigma}{2},
 \qquad
 y_3=\frac{-s-\sigma}{2}.
 \tag{F11}\label{eq:quartic-second-pair}
\]
The four roots of the original quartic are
\[
 x_i=y_i-\frac A4,
 \qquad i=0,1,2,3.
 \tag{F12}\label{eq:quartic-roots}
\]

\begin{theorem}[Exact Ferrari factorization]
\label{thm:ferrari-factorization}
Assume~\eqref{eq:ferrari-parameters} and~\eqref{eq:rho-sigma}.  Then
\[
 Y^4+pY^2+qY+r
 =\prod_{i=0}^{3}(Y-y_i),
 \tag{F13}\label{eq:quartic-complete-factor}
\]
with $y_i$ as in~\eqref{eq:quartic-first-pair} and
\eqref{eq:quartic-second-pair}.  After the translation~\eqref{eq:quartic-depress},
\[
 f(X)=a\prod_{i=0}^{3}(X-x_i).
 \tag{F14}\label{eq:quartic-original-factor}
\]
\end{theorem}

\begin{proof}
The equations~\eqref{eq:ferrari-parameters} give the two-quadratic identity
\eqref{eq:ferrari-target}.  By the quadratic factorization theorem,
\[
 Y^2-sY+(m-t)=(Y-y_0)(Y-y_1)
\]
because its discriminant is $s^2-4(m-t)=\rho^2$, and
\[
 Y^2+sY+(m+t)=(Y-y_2)(Y-y_3)
\]
because its discriminant is $s^2-4(m+t)=\sigma^2$.  Multiplying these
identities proves~\eqref{eq:quartic-complete-factor}.  Substituting
$Y=X+A/4$ and multiplying by $a$ gives
\eqref{eq:quartic-original-factor}.
\end{proof}

\subsection{Completeness and multiplicity}

In the $q=0$ branch, equation~\eqref{eq:biquadratic-factor} is already the
complete four-factor certificate.  In the $q\ne0$ branch, any root of the
depressed quartic makes the product of the two quadratic factors zero.  The
field no-zero-divisors law reduces the problem to one of the two quadratic
equations; Theorem~\ref{thm:quadratic-factorization} then yields one of the
four values.  Repeated roots cause no difficulty: the list may contain equal
entries, while the factorization retains their multiplicities.

\subsection{Primitive-recursive verification}

The certificate checker first verifies the translation
\eqref{eq:quartic-depress}.  It then reads the supplied branch certificate:
$q=0$ selects the biquadratic identity; $q\ne0$ selects Cardano's cubic
resolver and Ferrari's equations.  In the latter branch it checks
\[
 s^2-(2m-p)=0,
 \quad 2st+q=0,
 \quad t^2-(m^2-r)=0,
\]
then expands~\eqref{eq:ferrari-target} and applies the already verified
quadratic certificates.  This composition is primitive recursive.  The
existing Lean file \code{Quartic.lean} proves precisely the parametric Ferrari
identity and its soundness/exhaustiveness once $m,s,t,\rho,\sigma$ are
supplied; the additional biquadratic branch is a direct composition of its
quadratic layer.
